\newpage
\subsection{Widerstandsmessung II: Statistische Auswertung}\label{sec:01-2}

\begin{highlightblock}[Wiederholte Messungen mit statistischen Abweichungen]
\begin{enumerate}[label=\textbf{\alph*:}, ref=\alph*]
	\item Statistisch unabhängige Messung:\label{itm:P2-1}
	\begin{itemize}
		\item Unterschiedliche Werte von Strom- / Spannungswerte + Messbereiche
		\item Messgeräte tauschen.
		\item Kabel ein- und ausstecken
	\end{itemize}
	\item Messreihe statistisch auswerten (Schätzwert des Mittelwertes $\bar{x}$ und Standardabweichung~$s$)
	\item Konfidenzintervalle für Vertrauensniveaus: $1-\alpha =\{68.26\%, 99.5\%\}$ 
	\item Diskussion des Einflusses der unterschiedlichen Spannungsniveaus.
\end{enumerate} 
\end{highlightblock}

\subsubsection{Verwendetes Material}

\begin{table}[h]
	\centering
	\caption{Verwendetes Material}
	\begin{tabularx}{\textwidth}{XXX}
		\textbf{Material} & \textbf{Verwendung} & \textbf{Anmerkungen}\\ \midrule
		&  &  \\
		&  &  \\
		&  &  \\
	\end{tabularx}	
\end{table}

\subsubsection{Versuchsaufbau}

\marginref{\ref{itm:P2-1}}
Hier wurde der gleiche Versuchsaufbau wie in Abschnitt~\ref{sec:01-1} verwendet und folgende Messwerte Aufgenommen:

\begin{table}[h]
	\centering
	\caption{Messwerte Widerstandsmessung II}\label{tab:01-2-values}
	\pgfplotstabletypeset[
		col sep=comma,
		header=true,
		every head row/.style={after row=\midrule},
		begin table={\begin{tabularx}{\textwidth}},
		end table={\end{tabularx}},
		columns/desc/.style={column name={Bemerkung}, column type=|X, string type},
		columns/Nr/.style={column name={Nr}, column type=c},
		columns/U/.style={column name={$U / \si{\volt}$}, column type=c, precision=2, fixed},
		columns/UV/.style={column name={$U_V / \si{\volt}$}, column type=c, precision=2, fixed},
		columns/IA/.style={column name={$I_A / \si{\milli\ampere}$}, column type=c, precision=3, fixed},
		columns/R/.style={column name={$R / \si{\ohm}$}, column type=c, precision=2, fixed},
	]{../data/12_Values.csv}
\end{table}

\reversemarginpar\marginnote{\textbf{(\ref{itm:P1-2})}}[5pt]
Die Berechnung des Mittelwertes ($\bar{x}$) und der Standardabweichung~$s$~\cite[S.~7]{EMTPraktikum2025}:
\begin{align}
	& \bar{x} = \frac{1}{n} \sum_{i=1}^{n} x_i = \SI{67.94}{\ohm} \\
	& s = \sqrt{\frac{1}{n-1} \sum_{i=1}^n {(x_i - \bar{x})}^2} = \SI{0.08451}{\ohm} \approx \SI{0.08}{\ohm}
\end{align}

\reversemarginpar\marginnote{\textbf{(\ref{itm:P1-3})}}[5pt]
Das erweiterte Unsicherheitsintervall berechnet sich nach~\cite[S.~10, Tabelle~1.1]{EMTPraktikum2025} zu
\[
	u_{z} = t_{0} \frac{s}{\sqrt{n}} \, .
\]

Für $n=5$ ergibt sich:

\begin{itemize}
	\item 68,26\% Vertrauensniveau ($t_{0}/\sqrt{n} = 0.51$):
		\[
			u_{z} = 0.51 \cdot \SI{0.08}{\ohm} \approx \SI{0.041}{\ohm},
		\]
		also das Konfidenzintervall
		\[
			\bar{x} \pm u_{z} = (67.94 \pm 0.04) \, \si{\ohm}.
		\]

	\item 99,5\% Vertrauensniveau ($t_{0}/\sqrt{n} = 2.50$):
		\[
			u_{z} = 2.50 \cdot \SI{0.08}{\ohm} \approx \SI{0.175}{\ohm},
		\]
		also das Konfidenzintervall
		\[
			\bar{x} \pm u_{z} = (67.94 \pm 0.18) \, \si{\ohm}.
		\]
\end{itemize}

\reversemarginpar\marginnote{\textbf{(\ref{itm:P1-4})}}[5pt]
Einflusses der unterschiedlichen Spannungsniveaus:

Am Anfang des Strommessbereichs ist der relative Fehler am kleinsten (vgl. Messung~3).
Dies entspricht den Herstellerangaben zum Digitalmultimeter METEX~4600
(vgl. \cite[S.~143]{EMTPraktikum2025}), wonach der Fehler als Prozentwert des Anzeigewerts
angegeben wird. Bei kleinen Strömen (niedriges Spannungsniveau) ergibt sich daher die
geringste Abweichung, während diese bei höheren Spannungsniveaus zunimmt.

